\documentclass[11pt,a4paper]{article}
\usepackage[utf8]{inputenc}
\usepackage{geometry}
\usepackage{graphicx}
\usepackage{booktabs}
\usepackage{hyperref}
\usepackage{amsmath}
\usepackage{amssymb}
\usepackage{natbib}
\usepackage{enumitem}
\usepackage{xcolor}
\usepackage{fancyhdr}
\usepackage{titlesec}

% ================= 图像使用规范（供生成模型参考） =================
% 使用如下标准格式插入图片：
%
% \begin{figure}[htbp]
% \centering
% \includegraphics[width=0.8\linewidth]{path/to/image}
% \caption{Descriptive caption explaining the figure in the context of the literature.}
% \label{fig:unique_label}
% \end{figure}
%
% 要求：
% 1. 图片路径必须使用原始路径（不要修改或重写）
% 2. 每个图片必须有 \caption 和 \label
% 3. 在正文中使用如下方式引用：
%    Figure~\ref{fig:unique_label}
% 4. 图片应放置在相关内容附近（不要集中在文末）
% ============================================================

% 页面设置
\geometry{left=2.5cm,right=2.5cm,top=2.5cm,bottom=2.5cm}

% 页眉页脚
\pagestyle{fancy}
\fancyhf{}
\fancyhead[L]{\leftmark}
\fancyhead[R]{Literature Review}
\fancyfoot[C]{\thepage}
\renewcommand{\headrulewidth}{0.4pt}

% 标题格式
\titleformat{\section}{\Large\bfseries\color{blue!50!black}}{\thesection}{1em}{}
\titleformat{\subsection}{\large\bfseries\color{blue!30!black}}{\thesubsection}{1em}{}

% 超链接颜色
\hypersetup{
    colorlinks=true,
    linkcolor=blue!60!black,
    citecolor=green!50!black,
    urlcolor=blue!70!black
}

\begin{filecontents}{references.bib}
% 示例参考文献，请替换为实际引用
@article{example2024,
  title={Example Paper for Literature Review},
  author={Author, Example},
  journal={Journal of Examples},
  volume={1},
  number={1},
  pages={1--10},
  year={2024},
  publisher={Example Publisher}
}

@inproceedings{sample2023,
  title={Sample Conference Paper},
  author={Sample, Author and Test, Researcher},
  booktitle={Proceedings of the Example Conference},
  pages={100--110},
  year={2023}
}
\end{filecontents}

% ========== 标题信息 ==========
\title{\textbf{Literature Review: [Your Research Topic]}}
\author{Your Name \\
        Your Institution \\
        \texttt{your.email@example.com}}
\date{\today}

\begin{document}

\maketitle

% ========== 摘要 ==========
\begin{abstract}
\noindent
\textbf{Purpose:} Briefly describe the purpose and scope of this literature review. \\
\textbf{Methods:} Describe your search strategy and selection criteria. \\
\textbf{Findings:} Summarize the key findings from the reviewed literature. \\
\textbf{Conclusions:} State the main conclusions and implications.

\vspace{0.5em}
\noindent\textbf{Keywords:} keyword1, keyword2, keyword3, keyword4
\end{abstract}

\tableofcontents
\newpage

% ========== 1. 引言 ==========
\section{Introduction}
\label{sec:intro}

\subsection{Background}
Provide background information on the research topic. Explain why this topic is important and relevant to your field of study.

\subsection{Research Questions}
Clearly state the research questions or objectives of this literature review:
\begin{enumerate}
    \item What are the main research questions you aim to answer?
    \item What specific aspects of the topic will you focus on?
    \item What gaps in the literature are you trying to identify?
\end{enumerate}

\subsection{Scope and Limitations}
Define the scope of your review:
\begin{itemize}
    \item Time period covered (e.g., papers from 2010-2024)
    \item Types of sources (e.g., peer-reviewed journals, conference papers)
    \item Geographic or domain limitations
    \item Any exclusion criteria
\end{itemize}

% ========== 2. 研究方法 ==========
\section{Methodology}
\label{sec:method}

\subsection{Search Strategy}
Describe your systematic search approach:
\begin{itemize}
    \item Databases used (e.g., Google Scholar, IEEE Xplore, PubMed)
    \item Search terms and keywords
    \item Date of search
    \item Any snowballing or citation tracking methods
\end{itemize}

\subsection{Selection Criteria}
Explain your inclusion and exclusion criteria:
\begin{itemize}
    \item \textbf{Inclusion criteria:} What characteristics must papers have to be included?
    \item \textbf{Exclusion criteria:} What characteristics would disqualify a paper?
\end{itemize}

\subsection{Data Extraction}
Describe how you extracted and organized information from the selected papers.

% ========== 3. 文献分类框架 ==========
\section{Thematic Framework}
\label{sec:framework}

Present your taxonomy or classification scheme for organizing the literature. This could be:

\subsection{Chronological}
Organize by time periods or development phases.

\subsection{Methodological}
Group papers by research methods used (e.g., experimental, theoretical, survey-based).

\subsection{Thematic}
Categorize by themes, topics, or research questions.

% ========== 4. 文献综述主体 ==========
\section{Review of Literature}
\label{sec:review}

\subsection{Theme 1: [First Major Theme]}
\label{subsec:theme1}

Summarize the key papers in this theme. For each paper, discuss:
\begin{itemize}
    \item Research objectives and questions
    \item Methodology used
    \item Key findings and contributions
    \item Strengths and limitations
\end{itemize}

\textbf{Key Papers:}
\begin{itemize}
    \item \cite{example2024} - Brief description of contribution
    \item \cite{sample2023} - Brief description of contribution
\end{itemize}

\subsection{Theme 2: [Second Major Theme]}
\label{subsec:theme2}

Continue with the same structure for each major theme...

\subsection{Theme 3: [Third Major Theme]}
\label{subsec:theme3}

Add more subsections as needed to cover all relevant themes...

% ========== 5. 比较分析 ==========
\section{Comparative Analysis}
\label{sec:comparison}

\subsection{Cross-Theme Comparison}
Compare and contrast findings across different themes. Use tables for clarity:

\begin{table}[htbp]
\centering
\caption{Comparison of Key Approaches}
\label{tab:comparison}
\begin{tabular}{@{}llll@{}}
\toprule
\textbf{Approach} & \textbf{Advantages} & \textbf{Limitations} & \textbf{Best For} \\
\midrule
Method A & Advantage 1, 2 & Limitation 1, 2 & Scenario X \\
Method B & Advantage 3, 4 & Limitation 3, 4 & Scenario Y \\
Method C & Advantage 5, 6 & Limitation 5, 6 & Scenario Z \\
\bottomrule
\end{tabular}
\end{table}

\subsection{Trend Analysis}
Identify trends and patterns in the literature over time:
\begin{itemize}
    \item Emerging research directions
    \item Shifting paradigms or methodologies
    \item Increasing/decreasing interest in specific topics
\end{itemize}

% ========== 6. 研究空白 ==========
\section{Research Gaps}
\label{sec:gaps}

Based on your review, identify the main gaps in the existing literature:

\begin{enumerate}
    \item \textbf{Gap 1:} Description of the first identified gap
    \item \textbf{Gap 2:} Description of the second identified gap
    \item \textbf{Gap 3:} Description of additional gaps
\end{enumerate}

For each gap, explain:
\begin{itemize}
    \item Why it is a significant gap
    \item What opportunities it presents for future research
    \item Potential challenges in addressing it
\end{itemize}

% ========== 7. 讨论 ==========
\section{Discussion}
\label{sec:discussion}

\subsection{Synthesis of Findings}
Synthesize the main findings from your review. What are the overarching conclusions?

\subsection{Theoretical Implications}
Discuss the theoretical contributions and implications of the reviewed work.

\subsection{Practical Implications}
Discuss the practical applications and real-world implications.

% ========== 8. 结论 ==========
\section{Conclusion}
\label{sec:conclusion}

\subsection{Summary}
Provide a concise summary of the literature review, including:
\begin{itemize}
    \item Main themes identified
    \item Key findings and trends
    \item Major research gaps
\end{itemize}

\subsection{Future Directions}
Suggest directions for future research based on the identified gaps:
\begin{enumerate}
    \item High-priority research questions
    \item Methodological improvements needed
    \item Emerging areas to explore
\end{enumerate}

\subsection{Limitations of This Review}
Acknowledge limitations of your review process itself.

% ========== 参考文献 ==========
\newpage
\bibliographystyle{apalike}
\bibliography{references}

% ========== 附录 ==========
\newpage
\appendix
\section{Search Results Summary}
\label{app:search}

Include supplementary information such as:
\begin{itemize}
    \item PRISMA flow diagram
    \item Complete list of search terms
    \item Detailed inclusion/exclusion criteria
    \item List of excluded papers with reasons
\end{itemize}

\section{Data Extraction Table}
\label{app:data}

Include a table summarizing key information from all reviewed papers:

\begin{table}[htbp]
\centering
\caption{Summary of Reviewed Papers}
\label{tab:papers}
\begin{tabular}{@{}p{2cm}p{3cm}p{3cm}p{4cm}@{}}
\toprule
\textbf{Author (Year)} & \textbf{Method} & \textbf{Key Findings} & \textbf{Limitations} \\
\midrule
Example (2024) & Method A & Finding 1 & Limitation 1 \\
Sample (2023) & Method B & Finding 2 & Limitation 2 \\
\bottomrule
\end{tabular}
\end{table}

\end{document}
